% ===================================================================
%  TABULKA č. 9 — Hory a lázeňské město. — Les a lov.
%  Traduction 2026 d'apres texto/io, controlee sur texto/fr, texto/ru
%  et texto/uk. Consigne : voir texto/cs/10-tabulka-01.tex.
%
%  LE TABLEAU DE LA MONTAGNE DONNE LE MEME RESULTAT QUE DANS LES SIX
%  COLONNES PRECEDENTES : le relief garde ses mots. « Skála » le
%  rocher, « stráň » la pente, « rokle » le ravin, « bystřina » le
%  torrent, « propast » l'abime, « pramen » la source, « mýtina » la
%  clairiere, « houština » le fourre, « kmen » le tronc, « kůra »
%  l'ecorce. C'est le septieme temoignage.
%
%  MAIS LA CURE THERMALE EST INTERNATIONALE ICI AUSSI : « lázeňské
%  město », « sanatorium », « minerální voda », « lanovka »,
%  « viadukt », « kasino », « vila ». Le meme partage qu'en galicien
%  au meme tableau, et a quatre alineas de distance.
%
%  ET LA MINE, ELLE, EST TCHEQUE JUSQU'AU DERNIER MOT : « důl »,
%  « šachta », « štola », « ruda », « huť », « vysoká pec »,
%  « litina », « lom », « pískovec », « krumpáč ». C'est le seul
%  vocabulaire technique du volume qui ne doive rien au latin, et il
%  n'est pas la par hasard : les mines de Boheme ont donne leurs mots
%  a l'allemand et non l'inverse.
% ===================================================================

%%K t09-apar-1 apar
\VUsaut{4.0mm}
\VUtitre{15.0pt}{60}{TABULKA č. 9}
\VUsaut{3.0mm}

%%K t09-tit-1 sub
\VUcentre{12.0pt}{0}{\textit{První scéna.}}
\VUsaut{3.0mm}

%%K t09-tit-2 sub
\VUpk{11.4pt}{40}{Hory. --- Lázeňské město.}
\VUsaut{3.0mm}

%%K t09-c1-01-1 p
1. --- Jsem už osm dní v Pratzu. Nemohu vyjádřit veškerý obdiv, který jsem pocítil,
když jsem stál před podivuhodným \VUgras{pohořím} \textsuperscript{(1)} Pyrenejí,
jehož \VUgras{hřeben} \textsuperscript{(2)} se rýsuje na modrém nebi jako obrovská
girlanda.

%%K t09-c1-02-1 p
2. --- Nepředstavoval jsem si, že pyrenejské \VUgras{vrcholy} \textsuperscript{(3)}
dosahují tak veliké \VUgras{výšky,} a zůstal jsem ohromen před překrásnými
\VUgras{ledovci} \textsuperscript{(5)} a zasněženými \VUgras{štíty}
\textsuperscript{(4)}, po nichž se valí tak hrozivé \VUgras{laviny.}

%%K t09-tit-3 sub
\VUsaut{3.0mm}
\VUpk{10.2pt}{40}{Horská krajina.}
\VUsaut{3.0mm}

%%K t09-c1-03-1 p
3. --- Hned nazítří po svém příjezdu jsem šel ke starému vyhaslému \VUgras{vulkánu}
\textsuperscript{(6)}, který je blízko Pratzu. Na holých a obnažených okrajích
\VUgras{kráteru} je ještě dobře vidět stopy po \VUgras{lávě}, která tekla před několika
staletími po \VUgras{svahu hory} \textsuperscript{(7)}. Podařilo se mi najít na jedné
ze \VUgras{strání} několik úlomků té lávy.

%%K t09-c1-04-1 p
4. --- Pratz leží na hranici a \VUgras{pevnost} \textsuperscript{(8)}, která hájí
\VUgras{soutěsku} \textsuperscript{(27)} oddělující Francii od Španělska, zcela
připomíná ty staré hrady na břehu Rýna, které se zdají číhat na \VUgras{kořist} z
vrcholu své skály.

%%K t09-c1-05-1 p
5. --- Obrovský \VUgras{orel} \textsuperscript{(9)}, který má své hnízdo nedaleko
\VUgras{tvrze} \textsuperscript{(8)}, často poutá mou pozornost. Ten \VUgras{dravý
pták}, jehož \VUgras{zobák} je silný, přináší často svým mláďatům jehně nebo kůzle,
které drží svými \VUgras{drápy.} Tak veliká jsou jeho \VUgras{křídla,} že může dlouho
plachtit se svým těžkým břemenem.

%%K t09-tit-4 sub
\VUsaut{3.0mm}
\VUpk{10.2pt}{40}{Výletník.}
\VUsaut{3.0mm}

%%K t09-c1-06-1 p
6. --- Nejpříjemnější procházkou je tam výlet k vodopádu \textsuperscript{(10)} v
«Kau». \VUgras{Vodopád} padá ve vírech a potom mizí jako krásný \VUgras{potok} v
\VUgras{jedlovém lese} \textsuperscript{(11)} položeném západně od města. Vystoupili
jsme tím lesem až k \VUgras{meteorologické observatoři} \textsuperscript{(12)} a z
vrcholu \VUgras{vyhlídky} jsme zahlédli celé \VUgras{panorama} kraje.
\VUgras{Krajina} je podivuhodná a \VUgras{příroda} se zdá zemi rozdávat všechny
poklady, jimiž vládne.

%%K t09-c1-07-1 p
7. --- K sestupu jsme použili \VUgras{lanovku} \textsuperscript{(13)} a zastavili jsme
se na malé \VUgras{stanici} blízko \VUgras{zřícenin} \textsuperscript{(14)} starého
\VUgras{feudálního hradu} kdysi obývaného pány z Pratzu.

%%K t09-c1-08-1 p
8. --- Ubohý hrad! Co si asi myslí, když dole vidí koketní \VUgras{lázeňské město}
\textsuperscript{(15)}, které je dnes módními \VUgras{lázněmi}?

%%K t09-c1-08-2 p
\VUgras{Podnebí} je tak příjemné, \VUgras{minerální voda} tak účinná, že
\VUgras{léčba} užívaná v \VUgras{sanatoriu} \textsuperscript{(17)} je opravdovým
potěšením.

%%K t09-tit-5 sub
\VUsaut{3.0mm}
\VUpk{10.2pt}{40}{Příjemné lázeňské město.}
\VUsaut{3.0mm}

%%K t09-c1-09-1 p
9. --- Ostatně se udělalo vše, aby byl pobyt příjemný. Byl postaven velký
\VUgras{viadukt} \textsuperscript{(16)}, aby umožnil železnici; \VUgras{park}
\textsuperscript{(18)} \VUgras{lázeňského ústavu}, v němž se \VUgras{koncertuje}, je
upraven půvabným způsobem. Několik půvabných «\VUgras{vil}»
\textsuperscript{(19)} bylo postaveno kolem kasina \textsuperscript{(21)} a velmi dobře udržovaná
\VUgras{silnice} \textsuperscript{(22)} značně usnadňuje spojení.

%%K t09-c1-10-1 p
10. --- Prostý \VUgras{násep} \textsuperscript{(23)} odděluje \VUgras{cestu}
\textsuperscript{(20)} od vlastního \VUgras{údolí} \textsuperscript{(25)}, na jehož dně
leží půvabné malé \VUgras{jezero} \textsuperscript{(26)}, které živí čirý
\VUgras{potok} \textsuperscript{(24)}.

%%K t09-c1-11-1 p
11. --- Na druhé straně údolí se těží \VUgras{železný důl} \textsuperscript{(28)}.
Vícekrát jsem šel vidět \VUgras{horníky} sjíždět do \VUgras{šachty} a dokonce jsem
vstoupil do \VUgras{štoly}, v níž tráví svůj den těžíce \VUgras{rudu,} která obsahuje
\VUgras{kov.}

%%K t09-c1-12-1 p
12. --- Blízko dolu byla postavena významná \VUgras{huť,} aby se hned zužitkovalo
bohatství, které se z něho těží. \VUgras{Vysoká pec} \textsuperscript{(29)}, vytápěná
\VUgras{černým uhlím}, kterého je zde hojnost, umožňuje rychle a lacino získat
\VUgras{litinu.}

%%K t09-c1-13-1 p
13. --- Nedaleko dolu se kope \VUgras{lom} \textsuperscript{(30)}. Ani \VUgras{žulu,}
ani \VUgras{porfyr,} ba ani \VUgras{pískovec} starý \VUgras{lamač} neuseká svým
\VUgras{krumpáčem,} nýbrž prostý \VUgras{kvádr} pro stavbu domů.

%%K t09-c1-14-1 p
14. --- Jednou z nejkrásnějších ozdob toho kraje jsou \VUgras{jedle}
\textsuperscript{(31)}. Všude na dně \VUgras{rokle} \textsuperscript{(32)}, na břehu
\VUgras{bystřiny} \textsuperscript{(33)}, \VUgras{propasti} \textsuperscript{(34)} nebo
srázu se jejich tenké jehlice zelenají.

%%K t09-tit-6 sub
\VUsaut{3.0mm}
\VUpk{10.2pt}{40}{Chůze pěšky.}
\VUsaut{3.0mm}

%%K t09-c1-15-1 p
15. --- Využívám svých prázdnin k \VUgras{dlouhým} procházkám. \VUgras{Cesty}
\textsuperscript{(35)}, které se vinou po hoře, umožňují ujít, bez ohledu na
vzdálenost, několik \VUgras{kilometrů} a často několik \VUgras{desítek kilometrů.}

%%K t09-c1-15-2 p
\VUgras{Mezníky} \textsuperscript{(36)} a \VUgras{ukazatele}
\textsuperscript{(37)} značí cesty, na nichž se často potkává mladý \VUgras{horal}
\textsuperscript{(38)} s \VUgras{brašnou} \textsuperscript{(40)} po boku, nesoucí
\VUgras{svišť} \textsuperscript{(39)}, nebo nějaký \VUgras{cikán}
\textsuperscript{(41)}, jehož \VUgras{kožich} \textsuperscript{(42)} podivně
kontrastuje se starou roztrhanou \VUgras{čapkou} \textsuperscript{(43)}. Vede na
\VUgras{řetěze medvěda} \textsuperscript{(44)} cvičeného.

%%K t09-tit-7 sub
\VUpk{10.2pt}{40}{Výstup na hory.}
\VUsaut{3.0mm}

%%K t09-c1-16-1 p
16. --- Téměř každý den chodím s \VUgras{průvodcem} \textsuperscript{(48)} cvičit
\VUgras{výstup} a jsem velmi hrdý, když s \VUgras{batohem} \textsuperscript{(47)} na
zádech a s «\VUgras{alpenstockem}» v ruce, jako opravdový \VUgras{turista}
\textsuperscript{(46)}, útočím na \VUgras{skály} \textsuperscript{(45)}.

%%K t09-c1-17-1 p
17. --- Z vrcholu hory vypadají obyvatelé roviny ne větší než mravenci; \VUgras{stádo
ovcí} \textsuperscript{(50)} pasoucí se na \VUgras{pastvině} \textsuperscript{(49)} se
zdá být hračkou pro děti. \VUgras{Borovice} \textsuperscript{(51)} nebo i
\VUgras{dřevěný domek} \textsuperscript{(52)} vypadá sotva jako skvrna na zeleni.

%%K t09-c1-18-1 p
18. --- Při těch výletech často potkáváme na \VUgras{pěšince} \textsuperscript{(53)}
\VUgras{lovce kamzíků} \textsuperscript{(54)}, který nese na rameni zvíře, jež právě
zabil.

%%K t09-c1-19-1 p
19. --- Uložil jsem do svého herbáře velmi pozoruhodnou větev \VUgras{kapradí}
\textsuperscript{(55)} a krásnou růžovou větvičku \VUgras{vřesu}
\textsuperscript{(57)}, které jsem natrhal u vchodu do malé \VUgras{jeskyně}
\textsuperscript{(56)} při své poslední procházce.

%%K t09-tit-8 sub
\VUsaut{3.0mm}
\VUcentre{12.0pt}{0}{\textit{Druhá scéna.}}
\VUsaut{3.0mm}

%%K t09-tit-9 sub
\VUpk{11.4pt}{40}{Les. --- Lov.}
\VUsaut{3.0mm}

%%K t09-c2-01-1 p
1. --- Včera jsem strávil rozkošný den v lese. Píle, která vládne mezi zvířaty a
\VUgras{hmyzem} \textsuperscript{(5)}, jimiž je obydlen, mě velmi zajímala.

%%K t09-c2-01-2 p
\VUgras{Pavouk} \textsuperscript{(1)}, který tká svou \VUgras{síť}
\textsuperscript{(2)} mezi dvěma větvemi, aby chytil \VUgras{mouchu}
\textsuperscript{(3)}, jež poletuje kolem, \VUgras{hlemýžď} \textsuperscript{(4)},
který s námahou nese svou ulitu, roháč pátrající po kořisti, pilný mravenec velmi
horlivý blízko \VUgras{mraveniště} \textsuperscript{(6)}, ti všichni malí chodili,
přicházeli, pracovali s neobyčejnou pílí.

%%K t09-c2-02-1 p
2. --- \VUgras{Had} \textsuperscript{(7)}, kterého jsem na první pohled pokládal za
\VUgras{zmiji}, ale který nebyl nic jiného než prostá \VUgras{užovka}, se schoulil pod
kmenem stromu a chvílemi vystrkoval svou tenkou hlavičku, která se stále zdála být
připravena kousnout. Přede mnou tlustý \VUgras{slimák} \textsuperscript{(8)} těžce lezl
poté, co pozřel jednu z chutných \VUgras{jahůdek} \textsuperscript{(11)}, které tam
hojně rostou.

%%K t09-c2-03-1 p
3. --- Půda byla pokryta \VUgras{lesními květinami}: \VUgras{prvosenkami}
\textsuperscript{(9)}, \VUgras{brčály} \textsuperscript{(10)} a mnoha jinými
rostlinami, které jsem neznal, kvetly mezi \VUgras{trsy trávy}
\textsuperscript{(12)}.

%%K t09-c2-04-1 p
4. --- V tom kraji je \VUgras{lanýž} téměř tak běžný jako \VUgras{houba}
\textsuperscript{(14)}, a \VUgras{sele} \textsuperscript{(13)} patřící hajnému mlsně
pátralo, zda se mu podaří nějaké najít.

%%K t09-tit-10 sub
\VUsaut{3.0mm}
\VUpk{10.2pt}{40}{Pytláctví.}
\VUsaut{3.0mm}

%%K t09-c2-05-1 p
5. --- Byl jsem přítomen \VUgras{konci} výjevu, který mě velmi pobavil. S pomocí čichu
svého \VUgras{jezevčíka} \textsuperscript{(15)} právě \VUgras{hajný}
\textsuperscript{(16)} překvapil \VUgras{pytláka} \textsuperscript{(22)} ve chvíli, kdy
se ten chystal odnést \VUgras{zvěř} \textsuperscript{(17)}, kterou zabil. Pytlák dal
přednost tomu, aby opustil svou kořist, než aby se dal zatknout, a utekl, a
\VUgras{zajíc} \textsuperscript{(18)}, \VUgras{laň} \textsuperscript{(19)},
\VUgras{srnec} \textsuperscript{(20)} a \VUgras{divočák} \textsuperscript{(21)} zůstali
ležet na trávě k radosti hajného.

%%K t09-c2-06-1 p
6. --- Rozmanitost stromů, které les obsahuje, je udivující. \VUgras{Buky}
\textsuperscript{(23)} a \VUgras{břízy} \textsuperscript{(26)}, \VUgras{borovice},
které dávají \VUgras{pryskyřici} \textsuperscript{(27)}, \VUgras{duby}
\textsuperscript{(29)} a také mnohé jiné mísí své větvoví.

%%K t09-c2-06-2 p
\VUgras{Břečťan} \textsuperscript{(24)}, který přilne ke kmeni vysokých stromů, barví
\VUgras{hvozd} \textsuperscript{(25)} zářivou zelení, jež se příjemně spojuje se
světlejším a bledě zeleným odstínem \VUgras{mechu} \textsuperscript{(31)}, v němž
\VUgras{datel} \textsuperscript{(30)} hledá hmyz.

%%K t09-tit-11 sub
\VUsaut{3.0mm}
\VUpk{10.2pt}{40}{Lesní krajina.}
\VUsaut{3.0mm}

%%K t09-c2-07-1 p
7. --- Staré duby mě okouzlily. Jejich tlusté zkroucené \VUgras{kořeny}
\textsuperscript{(32)}, napůl vystupující ze země, jim dávají skvělý vzhled síly a
statnosti, a ptám se, jak se \VUgras{majitel} \textsuperscript{(35)} může odhodlat dát
je porazit a vydat je \VUgras{pile} \textsuperscript{(34)} \VUgras{dřevorubce}
\textsuperscript{(33)}. Ubohé \VUgras{kmeny} \textsuperscript{(36)}, obnažené od své
\VUgras{kůry} \textsuperscript{(37)} nebo rozštípnuté \VUgras{sekerou}
\textsuperscript{(38)} \VUgras{nádeníka} \textsuperscript{(39)}, mě rmoutí.

%%K t09-c2-08-1 p
8. --- Blízko starý \VUgras{uhlíř} \textsuperscript{(41)} připravil svou
\VUgras{lopatou hromadu} \textsuperscript{(42)} uhlí a stařena nesla \VUgras{otep}
\textsuperscript{(40)} \VUgras{mrtvého dřeva} z větviček poražených stromů.

%%K t09-tit-12 sub
\VUsaut{3.0mm}
\VUpk{10.2pt}{40}{Lov.}
\VUsaut{3.0mm}

%%K t09-c2-09-1 p
9. --- Chtěl jsem si jít na chvíli odpočinout do \VUgras{hajovny}
\textsuperscript{(46)}, ale když jsem přišel blízko \VUgras{jezírka}
\textsuperscript{(43)}, na němž plave krásná \VUgras{labuť} \textsuperscript{(45)},
všiml jsem si, že nedávná \VUgras{povodeň} \textsuperscript{(44)} proměnila cestu v
opravdovou \VUgras{bažinu}. Musel jsem se obrátit a, procházeje blízko \VUgras{cedru}
\textsuperscript{(49)}, \VUgras{topolů} \textsuperscript{(48)} a \VUgras{jilmu}
\textsuperscript{(47)}, které jsou na okraji lesa, jsem viděl vyrazit z
\VUgras{mlází} \textsuperscript{(54)} překrásného \VUgras{jelena}
\textsuperscript{(50)}, pronásledovaného vytrvalou \VUgras{smečkou}
\textsuperscript{(51)}, kterou pobízel i \VUgras{roh} \textsuperscript{(53)}
\VUgras{lovců na koních} \textsuperscript{(52)}. Ubohé zvíře bylo zcela vyčerpáno.
Padlo umírajíc několik kroků ode mne.

%%K t09-c2-10-1 p
10. --- Hajného syn, kterého přilákal hluk \VUgras{štvanice}, vyšel z domu a vraceli
jsme se dva do lesa. Uviděl \VUgras{straku} \textsuperscript{(56)} na větvi starého
dubu. Myslel si, že její hnízdo je asi skryto v \VUgras{listoví}
\textsuperscript{(58)}, a chtěl to hnízdo chytit.

%%K t09-c2-10-2 p
Hbitý jako \VUgras{veverka} \textsuperscript{(61)} šplhal z \VUgras{větve}
\textsuperscript{(55)} na větev, ale nenašel nic a podařilo se mu jen vyplašit starou
\VUgras{vránu} \textsuperscript{(60)} sedící na \VUgras{ratolesti}
\textsuperscript{(57)}.
